% block: complex-symmetric-systems-heading pages:102,103 type:section
\subsection{{复对称线性方程组}}\label{sec:complex-symmetric-linear-systems}
{本节考虑下列形式的复对称线性方程组}
\[A\bm x=\bm b,\]
{其中 $A$ 为复对称矩阵, 即 $A=A^{\top}$ 且 $A\ne A^{\mathrm H}$. 注意, 复对称矩阵 $A$ 不是 Hermitian 矩阵. 下面是复对称矩阵的一个例子:}
\[A=\begin{bmatrix}1+i&2+i\\2+i&3\end{bmatrix}.\]
{关于复对称线性方程组的应用, 见 Section\ \ref{subsec:helmholtz-difference} 中的 Helmholtz 方程和 Section\ \ref{subsec:electronic-structure} 中的大规模电子结构计算. 本节介绍求解复对称线性方程组的 COCG 方法, COCR 方法和 QMR\_SYM 方法. 若将这些方法应用于实对称线性方程组, 则它们分别退化为 CG 方法, CR 方法和 MINRES 方法. 关于复对称线性方程组的其他 Krylov 子空间方法及应用, 例如见 \BookCite{1}, \BookCite{30}, \BookCite{37}, \BookCite{40}, \BookCite{83}, \BookCite{84}, \BookCite{124}, \BookCite{141} 及其中的参考文献.}
% block: cocg-heading page:103 type:subsection
\subsubsection{{共轭正交共轭梯度 (COCG) 方法}}\label{subsec:cocg-method}
{COCG 方法 \BookCite{197} 是求解复对称线性方程组最著名的方法. 将点积 $(\bm a,\bm b)$ 改为 $(\bar{\bm a},\bm b)$, 可从 CG 方法形式上导出该算法. 这里 $\bar{\bm a}$ 表示 $\bm a$ 的复共轭. 经过这一改写, COCG 方法的算法见 Algorithm\ \ref{alg:cocg}.}
% block: cocg-algorithm page:103 type:ordinary-paragraph
\begin{algorithm}[H]\caption{{COCG 方法 (注意: $(\bar{\bm a},\bm b)=\bm a^{\top}\bm b$)}}\label{alg:cocg}\begin{algorithmic}[1]
\Statex\textbf{Input:} $\bm x_0\in\mathbb C^N,\ \beta_{-1}=0,\ \bm p_{-1}=\bm0,\ \bm r_0=\bm b-A\bm x_0$
\Statex\textbf{Output:} $\bm x_n$
\For{$n=0,1,2,\ldots$, {直至收敛}}
\State $\bm p_n=\bm r_n+\beta_{n-1}\bm p_{n-1}$
\State $\alpha_n=\frac{(\bar{\bm r}_n,\bm r_n)}{(\bar{\bm p}_n,A\bm p_n)}$
\State $\bm x_{n+1}=\bm x_n+\alpha_n\bm p_n$
\State $\bm r_{n+1}=\bm r_n-\alpha_nA\bm p_n$
\State $\beta_n=\frac{(\bar{\bm r}_{n+1},\bm r_{n+1})}{(\bar{\bm r}_n,\bm r_n)}$
\EndFor\end{algorithmic}\end{algorithm}
% block: cocg-orthogonality page:103 type:ordinary-paragraph
{COCG 方法的残差向量 $\bm r_n$ 满足下列关系:}
\[\bm r_n\perp\overline{\mathcal K_n(A,\bm r_0)},\]
{其中 $\overline{\mathcal K_n(A,\bm r_0)}:=\mathcal K_n(\bar A,\bar{\bm r}_0)$. 由 Algorithm\ \ref{alg:cocg} 容易看出, 若矩阵 $A$ 为实对称矩阵, 则 COCG 方法与 CG 方法相同. 预条件 COCG 方法见 Algorithm\ \ref{alg:pcocg}.}
% block: pcocg-algorithm page:104 type:ordinary-paragraph
\begin{algorithm}[H]\caption{{预条件 COCG 方法 (注意: $(\bar{\bm a},\bm b)=\bm a^{\top}\bm b$)}}\label{alg:pcocg}\begin{algorithmic}[1]
\Statex\textbf{Input:} $\bm x_0\in\mathbb C^N,\ \bm r_0=\bm b-A\bm x_0,\ \bm p_{-1}=\bm0,\ \beta_{-1}=0$
\Statex\textbf{Output:} $\bm x_n$
\For{$n=0,1,2,\ldots$, {直至收敛}}
\State $\bm p_n=K^{-1}\bm r_n+\beta_{n-1}\bm p_{n-1}$
\State $\alpha_n=\frac{(\overline{K^{-1}\bm r_n},\bm r_n)}{(\bar{\bm p}_n,A\bm p_n)}$
\State $\bm x_{n+1}=\bm x_n+\alpha_n\bm p_n$
\State $\bm r_{n+1}=\bm r_n-\alpha_nA\bm p_n$
\State $\beta_n=\frac{(\overline{K^{-1}\bm r_{n+1}},\bm r_{n+1})}{(\overline{K^{-1}\bm r_n},\bm r_n)}$
\EndFor\end{algorithmic}\end{algorithm}
% block: cocr-heading page:104 type:subsection
\subsubsection{{共轭正交共轭残差 (COCR) 方法}}\label{subsec:cocr-method}
{COCR 方法 \BookCite{171} 也是求解复对称线性方程组的 Krylov 子空间方法. 将点积 $(\bm a,\bm b)$ 改为 $(\bar{\bm a},\bm b)$, 可由 CR 方法形式上导出该算法. COCR 方法产生的近似解使残差向量满足}
\[\bm r_n\perp\overline{A\mathcal K_n(A,\bm r_0)},\]
{其中 $\overline{A\mathcal K_n(A,\bm r_0)}:=\bar A\mathcal K_n(\bar A,\bar{\bm r}_0)$. 由此, 若矩阵 $A$ 为实对称矩阵, COCR 方法就与 CR 方法相同. 这表明, 若 $A$ 复对称且接近实矩阵, 则可以预期残差的 2-范数趋向于几乎单调下降, 或表现出平滑的收敛行为. 下面推导 COCR 方法. 首先在 Algorithm\ \ref{alg:conjugate-a-orthogonalization} 中引入共轭 $A$-正交化过程.}
% block: conjugate-a-orthogonalization-algorithm page:104 type:ordinary-paragraph
\begin{algorithm}[H]\caption{{不作归一化的共轭 $A$-正交化过程 (注意: $(\bar{\bm a},\bm b)=\bm a^{\top}\bm b$)}}\label{alg:conjugate-a-orthogonalization}\begin{algorithmic}[1]
\Statex\textbf{Input:} $\bm v_0=\bm r_0=\bm b-A\bm x_0$
\Statex\textbf{Input:} $t_{0,0}=(\overline{A\bm v_0},A\bm v_0)/(\bar{\bm v}_0,A\bm v_0)$
\Statex\textbf{Input:} $\bm v_1=-\alpha_0(A\bm v_0-t_{0,0}\bm v_0)$
\State $t_{0,0}=(\overline{A\bm v_1}A\bm v_1)/(\bar{\bm v}_1,A\bm v_1)$
\State $\bm v_1=-\alpha_0(A\bm v_1-t_{0,0}\bm v_1)$
\For{$n=1,2,\ldots,N-1$}
\State $t_{n-1,n}=(\overline{A\bm v_{n-1}},A\bm v_n)/(\bar{\bm v}_{n-1},A\bm v_{n-1})$
\State $t_{n,n}=(\overline{A\bm v_n},A\bm v_n)/(\bar{\bm v}_n,A\bm v_n)$
\State $\bm v_{n+1}=-\alpha_n(A\bm v_n-t_{n,n}\bm v_n-t_{n-1,n}\bm v_{n-1})$
\EndFor\end{algorithmic}\end{algorithm}
% block: cocr-matrix-form pages:104,105 type:ordinary-paragraph
{与 Lanczos 过程 (Algorithm\ \ref{alg:lanczos}) 类似, Algorithm\ \ref{alg:conjugate-a-orthogonalization} 中的向量 $\bm v_i$ 满足下列关系:}
\[\operatorname{span}\{\bm v_1,\bm v_2,\ldots,\bm v_n\}=\mathcal K_n(A,\bm r_0),\]
\begin{equation}\label{eq:cocr-a-orthogonality}(\bar{\bm v}_i,A\bm v_j)=0\quad\text{for }i\ne j.\end{equation}
{下面从 Algorithm\ \ref{alg:conjugate-a-orthogonalization} 推导 COCR 方法. 对给定的初始猜测 $\bm x_0$, 令 $\bm r_0:=\bm b-A\bm x_0$ 为初始残差向量. 则 Algorithm\ \ref{alg:conjugate-a-orthogonalization} 可表示为下列矩阵形式:}
\[A\underbrace{[\bm r_0,\bm r_1,\ldots,\bm r_{n-1}]}_{R_n}
=\underbrace{[\bm r_0,\bm r_1,\ldots,\bm r_{n-1},\bm r_n]}_{R_{n+1}}
\underbrace{\begin{bmatrix}t_{0,0}&t_{0,1}&&\\-\alpha_0^{-1}&t_{1,1}&\ddots&\\&-\alpha_1^{-1}&\ddots&t_{n-2,n-1}\\&&\ddots&t_{n-1,n-1}\\&&&-\alpha_{n-1}^{-1}\end{bmatrix}}_{T_{n+1,n}},\]
{其中 $t_{k-1,k}=\frac{(\overline{A\bm r_{k-1}},A\bm r_k)}{(\bar{\bm r}_{k-1},A\bm r_{k-1})}$, $t_{k,k}=\frac{(\overline{A\bm r_k},A\bm r_k)}{(\bar{\bm r}_k,A\bm r_k)}$. 由上述矩阵形式得到}
\begin{equation}\label{eq:cocr-lanczos-matrix}AR_n=R_{n+1}T_{n+1,n}.\end{equation}
{由于矩阵 $R_n$ 由 Algorithm\ \ref{alg:conjugate-a-orthogonalization} 生成, 由 $A$-正交性质 \eqref{eq:cocr-a-orthogonality} 有 $R_n^{\top}AR_n=D_n$, 其中 $D_n$ 为 $n$ 行 $n$ 列的对角矩阵. 标量参数 $\alpha_0,\ldots,\alpha_{n-1}$ 仍未知. 我们将说明, 若确定这些标量参数, 使得能够从 $\bm r_0,\ldots,\bm r_{n-1}$ 的信息中提取近似解 $\bm x_0,\ldots,\bm x_{n-1}$, 就得到 COCR 方法的算法.}
% block: cocr-residual-extraction-condition pages:105,106 type:ordinary-paragraph
{设 $X_n:=[\bm x_0,\ldots,\bm x_{n-1}]$, $\bm1$ 为 $[1,\ldots,1]^{\top}$. 则 $R_n:=[\bm r_0,\bm r_1,\ldots,\bm r_{n-1}]$ 与 $X_n$ 有如下关系:}
\begin{equation}\label{eq:cocr-residual-solution-matrix}R_n=[\bm b-A\bm x_0,\bm b-A\bm x_1,\ldots,\bm b-A\bm x_{n-1}]=\bm b\bm1_n^{\top}-AX_n.\end{equation}
{将 \eqref{eq:cocr-residual-solution-matrix} 代入 \eqref{eq:cocr-lanczos-matrix}, 得}
\[AR_n=(\bm b\bm1_{n+1}^{\top}-AX_{n+1})T_{n+1,n}.\]
{于是,}
\[R_n=(\bm x\bm1_{n+1}^{\top}-X_{n+1})T_{n+1,n},\]
{其中 $\bm x$ 为线性方程组 $A\bm x=\bm b$ 的精确解. 若如下形式的条件}
\begin{equation}\label{eq:cocr-extraction-zero-condition}\bm x\bm1_{n+1}^{\top}T_{n+1,n}=O_n\end{equation}
{成立, 则 $R_n=-X_{n+1}T_{n+1,n}$. 关系 $R_n=-X_{n+1}T_{n+1,n}$ 表明, 可以由残差向量得到近似解向量. 反之, 若这一关系不成立, 就不能由残差向量得到近似解向量.}
{由于 $T_{n+1,n}$ 中的未知参数 $\alpha_0,\ldots,\alpha_{n-1}$ 尚未确定, 现确定这些参数使 \eqref{eq:cocr-extraction-zero-condition} 成立. 由此得到}
\[\begin{aligned}\alpha_0^{-1}&=t_{0,0},\\\alpha_k^{-1}&=t_{k,k}+t_{k-1,k},\quad1\le k\le n-1.\end{aligned}\]
% block: cocr-bidiagonal-factorization page:106 type:ordinary-paragraph
{将上述递推关系代入 $T_{n+1,n}$, 有}
\[T_{n+1,n}=\begin{bmatrix}\alpha_0^{-1}&t_{0,1}&&\\-\alpha_0^{-1}&\alpha_1^{-1}-t_{0,1}&\ddots&\\&-\alpha_1^{-1}&\ddots&t_{n-2,n-1}\\&&\ddots&\alpha_{n-1}^{-1}-t_{n-2,n-1}\\&&&-\alpha_{n-1}^{-1}\end{bmatrix}.\]
{于是, $T_{n+1,n}$ 可分解如下:}
\[\begin{aligned}
T_{n+1,n}&=\begin{bmatrix}1&&\\-1&\ddots&\\&\ddots&1\\&&-1\end{bmatrix}
\begin{bmatrix}\alpha_0^{-1}&&\\&\ddots&\\&&\alpha_{n-1}^{-1}\end{bmatrix}
\begin{bmatrix}1&\alpha_0t_{0,1}&&\\&1&\ddots&\\&&\ddots&\alpha_{n-2}t_{n-2,n-1}\\&&&1\end{bmatrix}\\
&=B_{n+1,n}^{(L)}\Omega_n^{-1}B_n^{(U)}.
\end{aligned}\]
{由 \eqref{eq:cocr-lanczos-matrix} 和上述分解得到}
\begin{equation}\label{eq:cocr-factorized-lanczos}
AR_n=R_{n+1}T_{n+1,n}=R_{n+1}B_{n+1,n}^{(L)}\Omega_n^{-1}B_n^{(U)}.
\end{equation}
{这里 $P_n:=R_n(B_n^{(U)})^{-1}$. 由 \eqref{eq:cocr-factorized-lanczos} 有}
\begin{equation}\label{eq:cocr-direction-matrix}
AP_n=AR_n(B_n^{(U)})^{-1}=R_{n+1}B_{n+1,n}^{(L)}\Omega_n^{-1}.
\end{equation}
{上述矩阵形式等价于下列递推关系:}
\begin{equation}\label{eq:cocr-residual-recurrence}\bm r_k=\bm r_{k-1}-\alpha_{k-1}A\bm p_{k-1},\quad1\le k\le n.\end{equation}
{由 $P_n=R_n(B_n^{(U)})^{-1}$ 得}
\begin{equation}\label{eq:cocr-direction-recurrence}\bm p_k=\bm r_k+\beta_{k-1}\bm p_{k-1},\quad1\le k\le n,\end{equation}
% block: cocr-solution-recurrence page:107 type:ordinary-paragraph
{其中 $\beta_{k-1}:=-\alpha_{k-1}t_{k-1,k}$. 现在给出近似解的计算公式. 由 \eqref{eq:cocr-direction-matrix},}
\[\begin{aligned}
AP_n&=R_{n+1}B_{n+1,n}^{(L)}\Omega_n^{-1}\\
&=(\bm b\bm1^{\top}-AX_{n+1})B_{n+1,n}^{(L)}\Omega_n^{-1}\\
&=-AX_{n+1}B_{n+1,n}^{(L)}\Omega_n^{-1}.
\end{aligned}\]
{因此得到}
\[P_n\Omega_n=X_{n+1}(-B_{n+1,n}^{(L)}),\]
{它等价于}
\begin{equation}\label{eq:cocr-solution-recurrence}\bm x_k=\bm x_{k-1}+\alpha_{k-1}\bm p_{k-1},\quad1\le k\le n.\end{equation}
% block: cocr-beta-and-bilinear-symmetry page:107 type:ordinary-paragraph
{现在给出 $\alpha_k$ 和 $\beta_k$ 更实用的计算公式. 由 \eqref{eq:cocr-lanczos-matrix} 得}
\[(AR_n)^{\top}AR_n=R_n^{\top}AR_{n+1}T_{n+1,n}.\]
{由上述关系的 $(k+1,k)$ 元素有 $(\overline{A\bm r_k},A\bm r_{k-1})=-\alpha_{k-1}^{-1}(\bar{\bm r}_k,A\bm r_k)$. 因此由 $\beta_{k-1}=-\alpha_{k-1}t_{k-1,k}$ 得}
\begin{equation}\label{eq:cocr-beta}
\beta_k=-\alpha_kt_{k,k+1}=\frac{(\bar{\bm r}_{k+1},A\bm r_{k+1})}{(\overline{A\bm r_{k+1}},A\bm r_k)}\times\frac{(\overline{A\bm r_k},A\bm r_{k+1})}{(\bar{\bm r}_k,A\bm r_k)}=\frac{(\bar{\bm r}_{k+1},A\bm r_{k+1})}{(\bar{\bm r}_k,A\bm r_k)}.
\end{equation}
{这里使用了关系 $(\bar{\bm r}_k,A\bm r_{k+1})=(\bar{\bm r}_{k+1},A\bm r_k)$, 因为}
\begin{equation}\label{eq:cocr-bilinear-symmetry}
(\bar{\bm r}_k,A\bm r_{k+1})=(\bar{\bm r}_k,A\bm r_{k+1})^{\top}=(\bm r_k^{\top}A\bm r_{k+1})^{\top}=(\bar{\bm r}_{k+1},A^{\top}\bm r_k)=(\bar{\bm r}_{k+1},A\bm r_k).
\end{equation}
% block: cocr-alpha-and-direction-orthogonality pages:107,108 type:ordinary-paragraph
{由 \eqref{eq:cocr-direction-matrix}, 并回顾 $P_n=R_n(B_n^{(U)})^{-1}$, 得}
\[(AP_n)^{\top}AP_n=(AP_n)^{\top}R_{n+1}B_{n+1,n}^{(L)}\Omega_n^{-1}
=(B_n^{(U)})^{-\top}R_n^{\top}AR_{n+1}B_{n+1,n}^{(L)}\Omega_n^{-1}.\]
{$(B_n^{(U)})^{-\top}$ 为下三角矩阵, $A$-正交性质 \eqref{eq:cocr-a-orthogonality} 表明 $R_n^{\top}AR_{n+1}$ 的非对角元素为零. 因而 \linebreak[4]$(AP_n)^{\top}AP_n$ 为如下形式的下三角矩阵}
\[(AP_n)^{\top}AP_n=\begin{bmatrix}d_0&&&\\ *&\ddots&&\\\vdots&\ddots&\ddots&\\ *&\cdots&*&d_{n-1}\end{bmatrix},\]
{其中 $d_k=(\bar{\bm r}_k,\bm r_k)\alpha_k^{-1}$. 因此得到}
\begin{equation}\label{eq:cocr-alpha}\alpha_k=\frac{(\bar{\bm r}_k,A\bm r_k)}{(\overline{A\bm p_k},A\bm p_k)}.\end{equation}
{此外, 由于 $(AP_n)^{\top}AP_n$ 是复对称的, 有}
\[(AP_n)^{\top}AP_n=\begin{bmatrix}d_0&&&\\ *&\ddots&&\\\vdots&\ddots&\ddots&\\ *&\cdots&*&d_{n-1}\end{bmatrix}
=\begin{bmatrix}d_0&*&\cdots&*\\&\ddots&\ddots&\vdots\\&&\ddots&*\\&&&d_{n-1}\end{bmatrix},\]
{这意味着 $(AP_n)^{\top}AP_n=\operatorname{diag}(d_0,\ldots,d_{n-1})$. 因此 $\bm p_i$ 和 $\bm p_j$ 共轭 $A^2$-正交, 即}
\[(\overline{A\bm p_i},A\bm p_j)=0\quad\text{for }i\ne j.\]
{由 \eqref{eq:cocr-residual-recurrence}--\eqref{eq:cocr-alpha} 得到 COCR 方法的算法, 见 Algorithm\ \ref{alg:cocr}.}
% block: cocr-algorithm page:108 type:ordinary-paragraph
\begin{algorithm}[H]\caption{{COCR 方法 (注意: $(\bar{\bm a},\bm b)=\bm a^{\top}\bm b$)}}\label{alg:cocr}\begin{algorithmic}[1]
\Statex\textbf{Input:} $\bm x_0\in\mathbb C^N,\ \beta_{-1}=0,\ \bm p_{-1}=\bm0,\ \bm r_0=\bm b-A\bm x_0$
\Statex\textbf{Output:} $\bm x_n$
\For{$n=0,1,2,\ldots$, {直至收敛}}
\State $\bm p_n=\bm r_n+\beta_{n-1}\bm p_{n-1}$
\State $A\bm p_n=A\bm r_n+\beta_{n-1}A\bm p_{n-1}$
\State $\alpha_n=\frac{(\bar{\bm r}_n,A\bm r_n)}{(\overline{A\bm p_n},A\bm p_n)}$
\State $\bm x_{n+1}=\bm x_n+\alpha_n\bm p_n$
\State $\bm r_{n+1}=\bm r_n-\alpha_nA\bm p_n$
\State $\beta_n=\frac{(\bar{\bm r}_{n+1},A\bm r_{n+1})}{(\bar{\bm r}_n,A\bm r_n)}$
\EndFor\end{algorithmic}\end{algorithm}
% block: pcocr-transformation page:108 type:ordinary-paragraph
{对线性方程组 $K_1^{-1}AK_1^{-\top}\widetilde{\bm x}=K_1^{-1}\bm b$ ($\widetilde{\bm x}=K_1^{\top}\bm x$) 应用 COCR 方法, 通过下列改写即可得到预条件 COCR 方法:}
\[\widetilde{\bm x}_n\Rightarrow K_1^{\top}\bm x_n,\quad\widetilde{\bm p}_n\Rightarrow K_1^{\top}\bm p_n,\quad\widetilde{\bm r}_n\Rightarrow K_1^{-1}\bm r_n,\]
{其中 $A\approx K=K_1K_1^{\top}$. 于是得到 Algorithm\ \ref{alg:pcocr} 所述的预条件 COCR 方法.}
% block: pcocr-algorithm page:109 type:ordinary-paragraph
\begin{algorithm}[H]\caption{{预条件 COCR 方法 (注意: $(\bar{\bm a},\bm b)=\bm a^{\top}\bm b$)}}\label{alg:pcocr}\begin{algorithmic}[1]
\Statex\textbf{Input:} $\bm x_0\in\mathbb C^N,\ \bm r_0=\bm b-A\bm x_0,\ \bm p_{-1}=\bm0,\ \beta_{-1}=0$
\Statex\textbf{Output:} $\bm x_n$
\For{$n=0,1,2,\ldots$, {直至收敛}}
\State $\bm p_n=K^{-1}\bm r_n+\beta_{n-1}\bm p_{n-1}$
\State $(A\bm p_n=AK^{-1}\bm r_n+\beta_{n-1}A\bm p_{n-1})$
\State $\alpha_n=\frac{(\overline{K^{-1}\bm r_n},AK^{-1}\bm r_n)}{(\overline{A\bm p_n},K^{-1}A\bm p_n)}$
\State $\bm x_{n+1}=\bm x_n+\alpha_n\bm p_n$
\State $\bm r_{n+1}=\bm r_n-\alpha_nA\bm p_n$
\State $(K^{-1}\bm r_{n+1}=K^{-1}\bm r_n-\alpha_nK^{-1}A\bm p_n)$
\State $\beta_n=\frac{(\overline{K^{-1}\bm r_{n+1}},AK^{-1}\bm r_{n+1})}{(\overline{K^{-1}\bm r_n},AK^{-1}\bm r_n)}$
\EndFor\end{algorithmic}\end{algorithm}
% block: pcocr-operation-count page:109 type:ordinary-paragraph
{注意, Algorithm\ \ref{alg:pcocr} 中加入第 3 和第 7 行, 是为了减少形如 $A\bm v$ 的矩阵向量乘法及形如 $K\bm v=\bm w$ 的线性方程组求解次数, 这与 Section\ \ref{subsec:conjugate-residual} 中预条件 CR 方法所述的技术相同.}
% block: qmr-sym-heading page:109 type:subsection
\subsubsection{{准最小残差 (QMR\_SYM) 方法}}\label{subsec:qmr-sym}
{van der Vorst 的著作 \BookCite{196}, 第 112 页中命名的 QMR\_SYM 方法 \BookCite{61}, 也是求解复对称线性方程组的一种著名 Krylov 子空间方法, 可从 Section\ \ref{subsec:complex-symmetric-lanczos} 所述的复对称 Lanczos 过程导出.}
{本小节介绍 QMR\_SYM 方法的推导过程. 设 $V_n$ 为 $N$ 行 $n$ 列的矩阵, 其列向量由复对称 Lanczos 过程 (Algorithm\ \ref{alg:complex-symmetric-lanczos}) 生成. 则 QMR\_SYM 方法从下式寻找近似解}
\begin{equation}\label{eq:qmr-sym-affine-solution}\bm x_n=\bm x_0+V_n\bm y_n,\quad\bm y_n\in\mathbb C^n.\end{equation}
{这里 $\bm y_n$ 由稍后所述的条件确定. 注意 $\bm x_n$ 位于仿射空间 $\bm x_0+\mathcal K_n(A,\bm r_0)$ 中. 相应的残差向量为}
\[\bm r_n=\bm r_0-AV_n\bm y_n.\]
{由复对称 Lanczos 过程的矩阵形式 \eqref{eq:complex-lanczos-matrix-relation} 得}
\[\bm r_n=\bm r_0-V_{n+1}T_{n+1,n}\bm y_n=V_{n+1}(\beta\bm e_1-T_{n+1,n}\bm y_n)\]
% block: qmr-sym-quasi-minimal-residual page:110 type:ordinary-paragraph
{其中 $\beta:=\|\bm r_0\|$. 上述推导过程与 Section\ \ref{subsec:minres} 中 MINRES 方法的推导相似. 区别在于, 若像 MINRES 方法一样选择 $\bm y_n$ 使残差范数达到极小值, 则可能导致大量计算开销, 因为}
\begin{equation}\label{eq:qmr-sym-residual-norm-distinction}
\|\bm r_n\|=\|V_{n+1}(\beta\bm e_1-T_{n+1,n}\bm y_n)\|\ne\|(\beta\bm e_1-T_{n+1,n}\bm y_n)\|,
\end{equation}
{这是由于 $V_{n+1}^{\mathrm H}V_{n+1}\ne I$. 这意味着, 若试图得到最小残差近似解, 就需要极小化 $\|V_{n+1}(\beta\bm e_1-T_{n+1,n}\bm y_n)\|$. 结果是计算开销和内存需求随迭代次数增加而增长.}
{因此, QMR\_SYM 方法不按最小残差途径设计, 而按准最小残差途径设计, 即选择 $\bm y_n$ 使得}
\begin{equation}\label{eq:qmr-sym-projected-minimization}\bm y_n:=\operatorname*{arg\,min}_{\bm y\in\mathbb C^n}\|\beta\bm e_1-T_{n+1,n}\bm y\|.\end{equation}
{所得向量 $\beta\bm e_1-T_{n+1,n}\bm y_n$ 称为准残差向量. 这里注意, 若 $V_n$ 满足关系 $V_n^{\mathrm H}V_n=I_n$, 则上述选择使残差范数达到极小值. 由于 $A$ 复对称, 除一些特殊情形外, 一般的 $V_n$ 不满足关系 $V_n^{\mathrm H}V_n=I_n$. 文献 \BookCite{61} 给出的一种特殊情形具有下列形式}
\[A=B+i\sigma I,\quad B=B^{\top}\in\mathbb R^{N\times N},\quad\sigma\in\mathbb R.\]
% block: qmr-sym-givens-solution page:110 type:ordinary-paragraph
{为实现极小化 \eqref{eq:qmr-sym-projected-minimization}, Section\ \ref{subsec:minres} 所述的 Givens 旋转起重要作用. 用满足 $Q_n^{\mathrm H}Q_n=I_n$ 和下式的 Givens 旋转 $Q_n=G_n\cdots G_1$ 乘以 $\beta\bm e_1-T_{n+1,n}\bm y$}
\[Q_nT_{n+1,n}=\begin{bmatrix}R_n\\\bm0^{\top}\end{bmatrix},\]
{有}
\begin{equation}\label{eq:qmr-sym-triangular-minimization}
\min_{\bm y\in\mathbb C^n}\|\beta\bm e_1-T_{n+1,n}\bm y\|=\min_{\bm y\in\mathbb C^n}\left\|\begin{bmatrix}\bm g_n\\g_{n+1}\end{bmatrix}-\begin{bmatrix}R_n\\\bm0^{\top}\end{bmatrix}\bm y\right\|,
\end{equation}
{其中}
\[\begin{bmatrix}\bm g_n\\g_{n+1}\end{bmatrix}=\beta Q_n\bm e_1\]
{且 $R_n$ 为上三角矩阵. 于是 \eqref{eq:qmr-sym-triangular-minimization} 的解为 $\bm y_n=R_n^{-1}\bm g_n$, 结合 \eqref{eq:qmr-sym-affine-solution}, 得到近似解 $\bm x_n=\bm x_0+V_nR_n^{-1}\bm g_n$. 与 MINRES 方法类似, 引入 $P_n:=V_nR_n^{-1}$ 可给出三项递推关系. 于是有 $\bm x_n$ 更实用的计算公式:}
% block: qmr-sym-residual-update page:111 type:ordinary-paragraph
\[\bm x_i=\bm x_{i-1}+g_i\bm p_i,\quad i=1,\ldots,n,\]
{其中 $g_i$ 和 $\bm p_i$ 分别为 $\bm g_n$ 的第 $i$ 个元素和 $P_n$ 的第 $i$ 列. QMR\_SYM 方法的完整算法见 Algorithm\ \ref{alg:qmr-sym}. 与 MINRES 方法不同, 由于关系 \eqref{eq:qmr-sym-residual-norm-distinction}, 不能从 \eqref{eq:qmr-sym-triangular-minimization} 得到残差范数 $\|\bm r_n\|$. 因此问题是如何计算 $\|\bm r_n\|$ 以评价近似解的质量. 一种简单办法是计算 $\|\bm b-A\bm x_n\|$, 但每次迭代需要额外一次矩阵向量乘法. 另一方面, 可以避免这次矩阵向量乘法计算. 下面给出这一技术. 与 MINRES 方法类似, QMR\_SYM 方法的近似解由下式给出}
\[\begin{aligned}\bm p_n&=\frac1{t_{n,n}}(\bm v_n-t_{n-2,n}\bm p_{n-2}-t_{n-1,n}\bm p_{n-1}),\\\bm x_n&=\bm x_{n-1}+g_n\bm p_n.\end{aligned}\]
{由于 $\bm r_n=\bm b-A\bm x_n$, 有}
\[\bm r_n=\bm r_{n-1}-g_nA\bm p_n.\]
{这里 $A\bm p_n$ 可按下式更新}
\[A\bm p_n=\frac1{t_{n,n}}(A\bm v_n-t_{n-2,n}A\bm p_{n-2}-t_{n-1,n}A\bm p_{n-1}).\]
{由于复对称 Lanczos 过程中已计算 $A\bm v_n$, 可以不直接计算 $\|\bm b-A\bm x_n\|$ 而计算残差的 2-范数 $\|\bm r_n\|$, 从而以较低开销评价 $\|\bm r_n\|$ 来检查收敛. 这些递推式加入 Algorithm\ \ref{alg:qmr-sym} 的第 24--26 行.}
% block: qmr-sym-algorithm page:112 type:ordinary-paragraph
\begin{algorithm}[H]\caption{{QMR\_SYM 方法 (注意: $(\bar{\bm a},\bm b)=\bm a^{\top}\bm b$)}}\label{alg:qmr-sym}\begin{algorithmic}[1]
\Statex\textbf{Input:} $\bm x_0\in\mathbb C^N,\ \beta_0=0,\ \bm v_0=\bm p_{-1}=\bm p_0=\bm0,\ \bm r_0=\bm b-A\bm x_0$
\Statex\textbf{Output:} $\bm x_n$
\State $\bm g=(\|\bm r_0\|,0,\ldots,0)^{\top},\ \bm v_1=\bm r_0/\|\bm r_0\|$
\For{$n=1,2,\ldots$}
\State ({复对称 Lanczos 过程})
\State $\alpha_n=(\bar{\bm v}_n,A\bm v_n)$
\State $\widetilde{\bm v}_{n+1}=A\bm v_n-\alpha_n\bm v_n-\beta_{n-1}\bm v_{n-1}$
\State $\beta_n=(\overline{\widetilde{\bm v}_{n+1}},\widetilde{\bm v}_{n+1})^{1/2}$
\State $\bm v_{n+1}=\widetilde{\bm v}_{n+1}/\beta_n$
\State $t_{n-1,n}=\beta_{n-1},\ t_{n,n}=\alpha_n$
\State $t_{n+1,n}=\beta_n$
\State (Givens {旋转})
\State \textbf{for} $i=\max\{1,n-2\},\ldots,n-1$ \textbf{do}
\State $\begin{bmatrix}t_{i,n}\\t_{i+1,n}\end{bmatrix}=\begin{bmatrix}c_i&s_i\\-\bar s_i&c_i\end{bmatrix}\begin{bmatrix}t_{i,n}\\t_{i+1,n}\end{bmatrix}$
\State \textbf{end for}
\State $c_n=\frac{|t_{n,n}|}{\sqrt{|t_{n,n}|^2+|t_{n+1,n}|^2}}$
\State $\bar s_n=\frac{t_{n+1,n}}{t_{n,n}}c_n$
\State $t_{n,n}=c_nt_{n,n}+s_nt_{n+1,n}$
\State $t_{n+1,n}=0$
\State $\begin{bmatrix}g_n\\g_{n+1}\end{bmatrix}=\begin{bmatrix}c_n&s_n\\-\bar s_n&c_n\end{bmatrix}\begin{bmatrix}g_n\\0\end{bmatrix}$
\State ({更新} $\bm x_n$)
\State $\bm p_n=(\bm v_n-t_{n-2,n}\bm p_{n-2}$
\State $\qquad-t_{n-1,n}\bm p_{n-1})/t_{n,n}$
\State $\bm x_n=\bm x_{n-1}+g_n\bm p_n$
\State ({检查收敛})
\State $A\bm p_n=(A\bm v_n-t_{n-2,n}A\bm p_{n-2}$
\State $\qquad-t_{n-1,n}A\bm p_{n-1})/t_{n,n}$
\State $\bm r_n=\bm r_{n-1}-g_nA\bm p_n$
\State {若 $\|\bm r_n\|/\|\bm b\|\le\epsilon$, 则停止}
\EndFor\end{algorithmic}\end{algorithm}
